\documentclass[10pt,a4paper]{article}
\usepackage[T1]{fontenc}
\usepackage[utf8]{inputenc}
\usepackage[brazilian]{babel}
\usepackage[a4paper,margin=1.5cm]{geometry}
\usepackage{booktabs}
\usepackage{xcolor}
\usepackage{microtype}
\usepackage[hidelinks]{hyperref}
\pagestyle{empty}
\setlength{\parindent}{0pt}
\setlength{\parskip}{2pt}
% Paleta monocromática: só a escala gray do xcolor, para impressão em toner.
\definecolor{cinza}{gray}{0.35}
\definecolor{cinzaclaro}{gray}{0.94}
\definecolor{borda}{gray}{0.75}
\definecolor{fundocard}{gray}{0.98}
\definecolor{barras}{gray}{0.15}
\newcommand{\barra}[1]{\textcolor{barras}{\rule{#1mm}{2.4pt}}}
% Cartão de indicador: rótulo, número em destaque e rodapé de contexto.
\newcommand{\kpicard}[3]{%
  \fcolorbox{borda}{fundocard}{%
    \begin{minipage}[c][1.4cm][c]{\dimexpr0.32\textwidth-2\fboxsep-2\fboxrule\relax}
      \centering
      {\scriptsize\bfseries\color{cinza} #1}\\[2pt]
      {\Large\bfseries #2}\\[1.5pt]
      {\tiny\color{cinza} #3}%
    \end{minipage}%
  }%
}
\begin{document}

{\Large\bfseries Barra do Chapéu --- climatologia das médias}\\[2pt]
{\small Estação CIIAGRO/IAC 293 \,\textbullet\,
99 meses medidos (2018-08 a 2026-10)}\\[1pt]
{\footnotesize \textbf{Atenção à série curta:} apenas 7 anos com os 12 meses fechados --- referência do período medido, não padrão consolidado.}

\vspace{4pt}
\colorbox{cinzaclaro}{%
\parbox{\dimexpr\textwidth-2\fboxsep}{%
\footnotesize
\textbf{Leia antes de usar:} os números abaixo são \textbf{médias de vários anos},
obtidas das medições mensais do CIIAGRO. \textbf{Não é previsão meteorológica} para
os próximos dias nem para esta safra.\par\vspace{2pt}
\textbf{Aplicação prática:} este é o comportamento médio esperado da estação ---
base técnica para o calendário agrícola, as janelas de calagem, o preparo do solo e
o manejo conservacionista.%
}}

\vspace{4pt}
\noindent
\kpicard{Chuva anual média}{1.261 mm/ano}{em todo o ano}\hfill
\kpicard{Mês mais chuvoso}{184 mm}{em janeiro}\hfill
\kpicard{Mês mais seco}{46 mm}{em julho}

\vspace{2pt}
\noindent
\kpicard{Temperatura média}{19,3 \textdegree{}C}{média dos 12 meses}\hfill
\kpicard{Mês mais quente}{22,8 \textdegree{}C}{em fevereiro}\hfill
\kpicard{Mês mais frio}{15,5 \textdegree{}C}{em julho}

\vspace{4pt}
\begin{center}
{\bfseries\small Tabela 1: Chuva média por mês do ano}\\[2pt]
{\small
\begin{tabular}{@{}l r l r r r@{}}
\toprule
mês & chuva (mm) & perfil & t mín & t méd & t máx \\
\midrule
janeiro    &   184 & \barra{26.0}   & 17,1 & 20,9 & 24,8 \\
fevereiro  &   132 & \barra{18.7}   & 17,0 & 22,8 & 28,5 \\
março      &   117 & \barra{16.5}   & 16,3 & 22,2 & 28,1 \\
abril      &    60 & \barra{8.5}    & 14,2 & 19,8 & 25,5 \\
maio       &    66 & \barra{9.3}    & 11,7 & 17,5 & 23,3 \\
junho      &    81 & \barra{11.5}   & 10,5 & 16,1 & 21,6 \\
julho      &    46 & \barra{6.5}    &  9,3 & 15,5 & 21,6 \\
agosto     &    67 & \barra{9.5}    & 10,3 & 16,2 & 22,1 \\
setembro   &   105 & \barra{14.8}   & 12,4 & 18,5 & 24,5 \\
outubro    &   140 & \barra{19.7}   & 13,9 & 18,9 & 24,0 \\
novembro   &   126 & \barra{17.8}   & 14,6 & 20,1 & 25,6 \\
dezembro   &   140 & \barra{19.7}   & 16,8 & 22,6 & 28,4 \\
\midrule
ano & 1.261 & & & 19,3 & \\
\bottomrule
\end{tabular}}
\end{center}

\begin{center}
{\bfseries\small Tabela 2: Chuva por estação do ano}\\[2pt]
{\small
\begin{tabular}{@{}l l r r r l l@{}}
\toprule
estação & meses & total (mm) & \% do ano & média/mês & maior mês & menor mês \\
\midrule
Verão      & dez--fev &   456 &  36\% &  152 & 184 (jan) & 132 (fev) \\
Outono     & mar--mai &   243 &  19\% &   81 & 117 (mar) & 60 (abr) \\
Inverno    & jun--ago &   195 &  15\% &   65 & 81 (jun) & 46 (jul) \\
Primavera  & set--nov &   371 &  29\% &  124 & 140 (out) & 105 (set) \\
\bottomrule
\end{tabular}}
\end{center}

\vspace{1pt}
{\footnotesize\textbf{Antes dos solstícios} (acumulado desde janeiro):
inverno: janeiro a maio somam 559 mm (44\% do ano) \,\textbullet\, verão: junho a novembro somam 566 mm (45\% do ano). Dezembro fica depois do solstício de verão e entra só no total
anual.}

\vspace{2pt}
{\bfseries\small Como usar no planejamento nesta estação}

\footnotesize\begin{itemize}\setlength{\itemsep}{0.5pt}
  \item \textbf{Amostragem de solo e preparo mecânico}: Em abril, maio, junho, julho e agosto a chuva média fica entre 46 e 81 mm/mês \textemdash{} é a janela de menor encharcamento para retirar terra para laudo, arar e incorporar calcário com antecedência.
  \item \textbf{Calagem em cobertura (pasto e perene já formado)}: A partir de outubro, quando a média volta a passar de 120 mm/mês: aí há água para o corretivo descer no perfil sem enxurrada.
  \item \textbf{Plantio e transplante}: O verão (dezembro a fevereiro) concentra 456 mm, 36\% da chuva do ano \textemdash{} favorece o pegamento de mudas sem irrigação pesada; em áreas declivosas, reforce a conservação de solo nesse período.
  \item \textbf{Calor, frio e fitossanidade}: fevereiro tem a maior média (22,8 \textdegree{}C) e julho a menor (15,5 \textdegree{}C, com mínima média de 9,3 \textdegree{}C). Antecipe o monitoramento de doenças fúngicas nos meses quentes e úmidos e proteja culturas sensíveis ao frio nos mais frios.
\end{itemize}

\vspace{1pt}
{\bfseries\small Vale para qualquer município da região}

\scriptsize\begin{itemize}\setlength{\itemsep}{0.5pt}
  \item \textbf{Média não é previsão}: Descrevem o padrão de vários anos, não o tempo deste ano: use para escolher a janela de cada operação, e decida no campo.
  \item \textbf{Corrija antes de precisar}: Calcário reage em meses; aplicado na janela seca, chega ao plantio já reagido.
  \item \textbf{Chuva em mm não é água no solo}: Textura, declive e cobertura morta decidem quanto do que cai chega à raiz.
  \item \textbf{Série curta é pista, não padrão}: Poucos anos fechados oscilam muito: compare com a estação vizinha antes de fixar calendário.
\end{itemize}

\vspace{1pt}
{\tiny\color{cinza}
Fonte: CIIAGRO/IAC (\href{http://ciiagro.org.br}{ciiagro.org.br}), série mensal
medida; chuva é o total do mês. Painel:
\href{https://climalitoralsul.netlify.app/\#barra-do-chapeu}%
{\texttt{climalitoralsul.netlify.app/\#barra-do-chapeu}}.
Umidade relativa omitida: a fonte usa \texttt{0} para ``sem sensor''.
médias do período 2018\textendash{}2026 (8.2 anos de dado).
Gerado a partir de \texttt{dados.json} de 2026-10-30.}

\end{document}
